\section{2021 Field Map: dónde estamos y qué falta}

Entender la arquitectura también exige entender por qué atrajo la atención de varios campos. Para 2021, el Transformer ya se había extendido de la generación y comprensión de lenguaje a vision, generative modeling, reinforcement learning y a sistemas relacionados con la biología estructural. Aquí no basta con enumerar nombres de aplicaciones; hay que preguntar cómo define cada campo el token, la position, la mask y el objective.

\subsection{La transferencia entre campos no es una copia simple}

En lenguaje, el token suele ser una subword; en imágenes, el token puede ser un patch; en RL se puede serializar la secuencia state/action/reward; las proteínas pueden modelarse con aminoácidos o con una residue representation. La misma fórmula de self-attention procesa relaciones completamente distintas en cada campo; por eso el preprocessing, el inductive bias y la evaluation determinan si la transferencia funciona.

\lecturefigure{slide-07-where-we-are.jpg}{En la clase de 2021 ya se ubicaba el Transformer en un mapa común de NLP, vision, RL y modelado generativo.}{Video oficial, 04:30.}

\begin{knowledgebox}{Lectura de la figura: primero identificar el token y después discutir la arquitectura}
En NLP, el token tiene un orden discreto de palabras y una gramática; un patch de vision además necesita posición bidimensional y estructura local; una trajectory de RL contiene la causalidad de las acciones y la recompensa; un generative image model puede trabajar sobre tokens de pixel, latent o codebook. Decir solo «todos usan Transformer» oculta las decisiones de modelado más importantes. Al leer las clases siguientes, conviene anotar primero cuatro columnas: token, position, mask, loss.
\end{knowledgebox}

\begin{warningbox}{AlphaFold2 no es «solo Transformer»}
La Transformer-like attention es un componente importante de AlphaFold2, pero el sistema también combina alineamientos de secuencias múltiples, pair representation, módulos geométricos e inferencia estructural. El éxito entre campos suele venir del diseño conjunto de la attention y la estructura propia del campo, no de aplicar sin cambios un modelo de texto a datos nuevos.
\end{warningbox}

\subsection{Los problemas persistentes en la Future slide}

Por un lado, la slide imagina aplicaciones como video understanding, finance y generative modeling; por otro, enumera los componentes que faltan: external memory, la reducción de la complejidad computacional y la alineación con los objetivos humanos. Estos problemas persisten porque el Transformer solo define cómo se mezcla la información; no aporta por sí mismo estado de largo plazo, cómputo de bajo costo ni objetivos correctos.

\lecturefigure{slide-08-future.jpg}{La expansión de aplicaciones y los componentes faltantes deben leerse a la vez.}{Video oficial, 05:55.}

\begin{knowledgebox}{Lectura de la figura: la columna derecha se parece más a una agenda de investigación que la izquierda}
Las aplicaciones de la izquierda responden «en qué tareas podría aportar valor»; los missing ingredients de la derecha responden «por qué la arquitectura por sí sola todavía no basta». La external memory se ocupa de la información persistente que excede el context; la efficiency, de la attention $O(n^2)$ y del costo de entrenamiento; la alignment, de los objetivos, la retroalimentación y las restricciones sociales. Toda proposal de aplicación debería corresponder al menos a una hipótesis de capacidad y a un modo de falla.
\end{knowledgebox}

\teachervoice{El ponente no describe el Transformer como un punto final, sino que coloca las aplicaciones junto a los componentes faltantes. Nota de clase: el criterio de investigación surge de ver al mismo tiempo «qué se puede hacer» y «por qué todavía no se hace bien», y no de dejar de preguntar después de una demo exitosa.}

\subsection{Resumen de la sección}

La unificación entre campos proviene de la interaction primitive, no de la homogeneización de las tareas. El mapa del futuro del curso deja la tokenización, la memory, la efficiency y la alignment como problemas pendientes, y ofrece así coordenadas para las guest lectures.
